\documentclass[11pt]{article}
\usepackage[margin=1in]{geometry}
\usepackage{amsmath,amssymb,amsthm}
\usepackage{graphicx}
\usepackage[hidelinks]{hyperref}

\newtheorem{theorem}{Theorem}
\newtheorem{lemma}{Lemma}
\newtheorem{remark}{Remark}

\newcommand{\E}{\mathbb E}
\newcommand{\R}{\mathbb R}
\newcommand{\ga}{\gamma}
\newcommand{\Lp}{L^p}

\title{The Smirnov Property for Every Nondegenerate Bivariate Gaussian Density}
\author{fa\_banach\_001, agent\_lane\_01}
\date{August 26, 2026}

\begin{document}
\maketitle

\section*{Status}

Candidate full solution, likely valid, pending human review.

\section*{Source Question}

Mayerhofer, \emph{The Smirnov Property for Weighted Lebesgue Spaces},
arXiv:2101.02470, Definition 2.4, calls a density $w$ Smirnov (for
$q>1$) when
\[
 F(x)+G(y)\in L^q(w),\qquad F,G\in L^1(w),
\]
forces $F,G\in L^q(w)$.  The harmless factor $1/2$ in the paper's
definition has been suppressed.  Remark 2.10 asks whether this property holds
for a bivariate standard normal density with nonzero correlation.

\begin{center}
\includegraphics[width=.97\linewidth]{figures/open_problem_crop.png}
\end{center}

We answer this question affirmatively.  In fact, the result holds for every
nondegenerate bivariate Gaussian law and every $1<q<\infty$.

\section*{Proof Idea}

After standardization, conditional expectation between the two Gaussian
coordinates is the one-dimensional Gaussian noise operator $P_\rho$.  If
$H=F(X)+G(Y)\in L^q$, then conditioning on $X$ and on $Y$ gives
\[
 A=F+P_\rho G\in L^q(\ga),\qquad
 B=P_\rho F+G\in L^q(\ga).
\]
The first summand cancels from $B-P_\rho A$, leaving
\[
 (I-P_{\rho^2})G\in L^q(\ga).
\]
The key rigidity fact is that an $L^1(\ga)$ function whose image under
$I-P_t$ lies in $L^q(\ga)$ must itself lie in $L^q(\ga)$, up to its constant
part.  On mean-zero functions, a sufficiently high power of $P_t$ is a strict
$L^q$ contraction.  This follows by interpolating its exact $L^2$ spectral
gap with the elementary $L^1$ or $L^\infty$ bound.  Thus $I-P_t$ is invertible
on mean-zero $L^q$.  An $L^1$ fixed-point argument rules out any residual
solution outside $L^q$.

\section*{Gaussian Noise Operators}

Let $\ga$ be standard Gaussian measure on $\R$.  For $-1<r<1$, define
\[
 (P_r f)(x)=\int_{\R}f\bigl(rx+\sqrt{1-r^2}\,z\bigr)\,d\ga(z).
\]
The operators $P_r$ are contractions on every $L^s(\ga)$,
$1\leq s\leq\infty$, by Jensen's inequality and invariance of $\ga$.  They
satisfy
\[
 P_rP_s=P_{rs},\qquad P_0=\Pi,
\]
where $\Pi f=\int f\,d\ga$ is the projection onto constants.  These identities
follow directly by combining the two independent Gaussian noises.

\begin{lemma}[Resolvent rigidity]\label{lem:rigidity}
Let $0\leq t<1$ and $1<q<\infty$.  If $g\in L^1(\ga)$ and
$(I-P_t)g\in L^q(\ga)$, then $g-\int g\,d\ga\in L^q(\ga)$.
\end{lemma}

\begin{proof}
Put $T=P_t-\Pi$, considered on the full space $L^s(\ga)$.  Since
$P_t\Pi=\Pi P_t=\Pi$, for every integer $m\geq1$,
\[
 T^m=P_{t^m}-\Pi.
\]
Consequently $\|T^m\|_{L^1\to L^1}\leq2$ and
$\|T^m\|_{L^\infty\to L^\infty}\leq2$.

On $L^2(\ga)$, the normalized Hermite polynomials form an orthonormal basis,
$P_sH_k=s^kH_k$, and $\Pi$ removes $H_0$.  Hence
\[
 \|T^m\|_{L^2\to L^2}=t^m.
\]
For completeness, the eigenvalue identity follows immediately by applying
$P_s$ to the generating function
$\exp(ux-u^2/2)=\sum_{k\geq0}H_k(x)u^k/k!$ and integrating the independent
Gaussian noise.

If $1<q<2$, interpolation between $L^1$ and $L^2$ gives
\[
 \|T^m\|_{L^q\to L^q}
 \leq 2^{1-\theta}t^{m\theta},
 \qquad \theta=2\left(1-\frac1q\right)>0.
\]
If $2<q<\infty$, interpolation between $L^2$ and $L^\infty$ gives the same
form of estimate with $\theta=2/q>0$.  For $q=2$ the displayed $L^2$ norm is
already enough.  Thus, in every case, some $m$ satisfies
$\|T^m\|_{L^q\to L^q}<1$.  Therefore $I-T^m$ is invertible on $L^q$ by a
Neumann series, and so is $I-T$, since
\[
 (I-T)^{-1}=(I+T+\cdots+T^{m-1})(I-T^m)^{-1}.
\]

Set $u=(I-P_t)g$.  Then $u\in L^q(\ga)$ and $\Pi u=0$.  Let
$v=(I-T)^{-1}u\in L^q(\ga)$.  Taking expectations in $(I-T)v=u$ shows
$\Pi v=0$; hence $Tv=P_tv$ and
\[
 (I-P_t)v=u.
\]
Now
\[
 d=g-\Pi g-v\in L^1(\ga),\qquad \Pi d=0,\qquad P_td=d.
\]
Iteration gives $d=P_{t^m}d$ for all $m$.  As $t^m\to0$,
$P_{t^m}d\to\Pi d=0$ in $L^1(\ga)$.  To see this last convergence without
any hidden regularity assumption, approximate $d$ in $L^1(\ga)$ by a bounded
continuous function $\phi$.  Dominated convergence gives
$P_s\phi\to\Pi\phi$ in $L^1(\ga)$ as $s\to0$, while the $L^1$ contraction
of $P_s$ controls the two approximation errors uniformly in $s$.
Thus $d=0$, and $g-\Pi g=v\in L^q(\ga)$.
\end{proof}

\section*{Main Result}

\begin{theorem}\label{thm:main}
Let $(X,Y)$ have a nondegenerate bivariate Gaussian distribution, and let
$1<q<\infty$.  If measurable functions $F$ and $G$ satisfy
\[
 F(X),G(Y)\in L^1,\qquad F(X)+G(Y)\in L^q,
\]
then $F(X),G(Y)\in L^q$.  Equivalently, every nondegenerate bivariate Gaussian
density has the Smirnov property for every $1<q<\infty$.
\end{theorem}

\begin{proof}
Separate affine changes of the two coordinates preserve all the stated
integrability properties.  We may therefore assume that $X$ and $Y$ are
standard normal with correlation $\rho$.  Nondegeneracy is precisely
$|\rho|<1$.

Write $H=F(X)+G(Y)$.  Conditional Jensen gives
\[
 A(X):=\E[H\mid X]\in L^q,
 \qquad
 B(Y):=\E[H\mid Y]\in L^q.
\]
The conditional laws of a standard bivariate Gaussian pair yield, as
one-variable identities in $L^1(\ga)$,
\[
 A=F+P_\rho G,\qquad B=P_\rho F+G.
\]
Since $P_\rho$ is an $L^q$ contraction and
$P_\rho^2=P_{\rho^2}$,
\[
 B-P_\rho A=(I-P_{\rho^2})G\in L^q(\ga).
\]
Lemma~\ref{lem:rigidity}, with $t=\rho^2<1$, implies
$G-\int G\,d\ga\in L^q(\ga)$, and constants belong to $L^q(\ga)$.
Thus $G\in L^q(\ga)$.  Finally,
\[
 F=A-P_\rho G\in L^q(\ga).
\]
This proves the claim.  (The same proof includes the independent case
$\rho=0$.)
\end{proof}

\section*{Consequence for the Source Integral Equations}

The remaining weight hypothesis in Mayerhofer's Theorem 3.1 can be evaluated
explicitly.  Let $p>1$ be the exponent in that theorem (so its Smirnov
exponent is the conjugate $q=p/(p-1)$).  For the standardized density
\[
 w_\rho(x,y)=\frac{1}{2\pi\sqrt{1-\rho^2}}
 \exp\!\left(-\frac{x^2-2\rho xy+y^2}{2(1-\rho^2)}\right),
\]
condition (2.3) asks that $w_Xw_Y/w_\rho$ belong to $L^p(w_\rho)$.  Apart
from a positive constant, the integrand in its $p$th power is the exponential
of minus one half of the quadratic form with matrix
\[
 M=pI-\frac{p-1}{1-\rho^2}
 \begin{pmatrix}1&-\rho\\-\rho&1\end{pmatrix}.
\]
The smaller eigenvalue of $M$ is
\[
 p-\frac{p-1}{1-|\rho|}
   =\frac{1-p|\rho|}{1-|\rho|}.
\]
Thus condition (2.3) holds exactly when
\[
 p|\rho|<1,
 \qquad\text{equivalently}\qquad
 q>\frac{1}{1-|\rho|}.
\]
Consequently, Theorem~\ref{thm:main} supplies the previously missing Smirnov
hypothesis, and Mayerhofer's uniqueness conclusion applies in this sharp
parameter range whenever the theorem's stated assumptions on the marginals
also hold.  Outside this range our Smirnov result remains valid, but the
source theorem cannot be invoked because its separate condition (2.3) fails.

\section*{Adversarial Checks}

The proof does not assume that $F$ or $G$ is initially square-integrable:
Hermite expansions are used only to estimate the bounded operator $T^m$ on
$L^2$, followed by interpolation.  The actual functions enter the spectral
argument only after $(I-T)^{-1}$ has been constructed on $L^q$.

The possible constant obstruction is handled explicitly by $\Pi$.  The
kernel of $I-P_t$ on $L^1(\ga)$ consists only of constants: the final
$L^1$ convergence argument proves this rather than importing an $L^2$
ergodicity claim.  Negative correlations cause no sign issue because the
composition used in the proof is $P_\rho^2=P_{\rho^2}$.

\section*{Novelty Check and Limitations}

The four local run indexes were searched for arXiv:2101.02470, the exact
source title, the bivariate-normal question, the Smirnov property, and
Gaussian/Ornstein--Uhlenbeck variants; no matching solution or attempt was
found.  The source context confirms that Remark 2.10 and the conclusion state
a genuine open problem.

A bounded web search on August 26, 2026 used the exact phrase ``Smirnov
property'' together with ``bivariate normal'', ``Gaussian'', and
``Ornstein--Uhlenbeck'', as well as the functional formulation involving
$F(X)+G(Y)$.  It found the arXiv paper and its 2024 published version, which
still describes the question as unknown, but no later resolution or this
noise-operator proof.  This was not a comprehensive MathSciNet or zbMATH
search, so the novelty assessment remains provisional.

The theorem treats exactly two Gaussian coordinates.  Extending the argument
to arbitrary multivariate Gaussian vectors would require controlling a coupled
system of conditional-expectation operators and is not claimed here.

\section*{Human Review Recommendation}

Send to a specialist in Gaussian analysis or weighted $L^p$ spaces.  The main
review points are the $L^q$ resolvent lemma, especially its $L^1$ fixed-point
closure, and the precise scope of the uniqueness corollary relative to all the
hypotheses of the source theorem.

\begin{thebibliography}{9}
\bibitem{Mayerhofer}
E. Mayerhofer,
\emph{The Smirnov Property for Weighted Lebesgue Spaces},
arXiv:2101.02470 (2021; revised 2024), published in
\emph{Mathematics} 12 (2024), no.~19, 3135,
\href{https://doi.org/10.3390/math12193135}{doi:10.3390/math12193135}.
\end{thebibliography}

\end{document}
